\section{Introduction}\label{sec:introduction}

Internet-of-Things (IoT) applications are increasingly built as collections
of microservices distributed across edge, fog, and cloud resources. This
architecture allows sensing, data preparation, inference, storage, and
control functions to scale and evolve independently. It also enables
latency-sensitive processing to remain close to devices while
compute-intensive services use elastic cloud capacity. The same flexibility,
however, makes scheduling considerably more difficult. An edge--cloud
scheduler must choose among heterogeneous nodes that differ not only in
available CPU and memory, but also in administrative domain, geographical
jurisdiction, trust level, encryption support, and hardware-isolation
capability.

These differences are particularly important for services that process
sensitive, regulated, or operationally critical data. For such services,
placement cannot be reduced to finding the node with the lowest load or
highest expected performance. A candidate node may offer excellent
performance but violate a data-residency rule. Another node may satisfy the
required jurisdiction but lack the necessary trust or isolation capability.
A third node may satisfy all policy requirements while exhibiting a rapidly
increasing workload that makes it a poor operational choice. The scheduling
decision must therefore answer several distinct questions: whether a node is
allowed to host the service, whether it is likely to remain operationally
suitable, whether sufficient headroom exists at commitment time, and which
eligible destination should be selected without repeatedly concentrating
decisions on the same small subset of nodes.

Most existing microservice schedulers primarily address the performance side
of this problem. Learning-based frameworks coordinate request dispatch and
service orchestration across heterogeneous edge--cloud infrastructures
\citep{shen2023kais}. Network-aware schedulers reduce communication overhead
by considering service dependencies, traffic intensity, or inter-node delay
\citep{chen2026trade}. Telemetry-driven controllers localize resource
contention, respond to service-level objective violations, and adapt resource
allocation or replica counts under changing workloads
\citep{qiu2020firm,gan2021sage,wen2024statuscale}. These approaches provide
important mechanisms for improving throughput, latency, and resource
efficiency, but their primary objective remains operational performance.

A parallel line of research focuses on whether distributed services and
worker nodes can be trusted. Policy-as-Code frameworks express and enforce
security requirements during placement, scaling, and migration
\citep{pallewatta2024policy}. Remote-attestation mechanisms establish a
hardware-rooted basis for admitting edge devices into cloud-native clusters
and for dynamically controlling their privileges
\citep{thijsman2024attested}. These mechanisms strengthen security and
governance, but they do not by themselves determine which legal node is the
best operational destination under rapidly changing resource conditions.

The separation between these two lines of work leaves an important practical
gap. Performance-oriented scheduling may rank a prohibited node highly,
whereas policy enforcement may identify the legal set without distinguishing
between stable and deteriorating destinations. Treating policy requirements
as finite penalties inside a weighted objective does not fully resolve this
gap. When legality is represented as another cost term, a sufficiently
attractive performance score may compensate for the penalty. Similarly,
ranking all nodes first and filtering prohibited destinations afterwards
mixes candidates that can be committed with candidates that should never have
entered the decision process.

For sensitive microservices, legality must instead define the feasible
decision space. A node that violates a mandatory jurisdiction, trust,
encryption, or isolation requirement should be excluded before any
performance comparison is performed. This legality-first principle gives
policy constraints a different role from optimization objectives: policy
determines which decisions are admissible, while optimization determines
which admissible decision is preferable.

Even within the legal candidate set, current utilization alone provides an
incomplete view of operational suitability. Two nodes can report the same CPU
and memory utilization at the decision instant while following very different
trajectories. A stable node and a node whose utilization has increased
sharply over recent intervals should not receive the same risk assessment.
Scheduling from a single telemetry snapshot therefore reacts only after
pressure becomes visible at the current instant. Short temporal histories,
including recent levels, changes, ranges, and variability, provide the
scheduler with information about the direction and persistence of the node
state.

Prediction alone is also insufficient. A risk model estimates the likelihood
of a near-term threshold crossing, but the selected destination must still
satisfy a transparent resource-headroom check immediately before commitment.
The two mechanisms serve complementary purposes: temporal prediction
distinguishes stable candidates from deteriorating ones, whereas a
pre-commit guardrail prevents a destination with insufficient projected
headroom from being selected. Keeping these roles separate avoids relying on
a single opaque score for both forecasting and executability.

Another challenge arises when multiple scheduling decisions are processed
within the same control interval. Repeatedly choosing the node with the
smallest predicted risk can concentrate placements on a few apparently
attractive destinations. Such concentration reduces placement diversity and
can create the next resource hotspot. Evaluating a large global fairness
objective for every event would introduce unnecessary overhead. A more
practical approach is to use historical selection counts to construct a
compact, fairness-aware shortlist and then apply operational ranking only
within that shortlist.

This paper presents \cage, a Compliance-Aware Guardrail-Enforced Scheduler
for sensitive edge--cloud microservices. \cage organizes scheduling as a
sequence of non-interchangeable decisions. First, a deterministic policy gate
constructs the legal candidate set from the service requirements and node
capabilities. Second, a calibrated temporal-risk model removes legal nodes
that are likely to cross the operational utilization boundary in the next
control interval. Third, a fairness-aware shortlist gives priority to
destinations that have not been repeatedly selected while retaining
operationally strong candidates. Finally, an empirical Q95 projection verifies
CPU and memory headroom before the destination is committed.

The central idea of \cage is therefore not to add policy, risk, fairness, and
capacity as competing terms in a single weighted objective. Instead, the
scheduler assigns each concern a precise role and places it at the appropriate
stage of the decision path. Mandatory policy requirements define
admissibility; temporal telemetry estimates near-term operational risk;
selection history controls concentration; and the pre-commit projection
checks immediate executability. This ordered design makes each decision
interpretable and prevents a favourable score in one dimension from
compensating for a mandatory failure in another.

The main contributions of this work are as follows.

\begin{enumerate}[leftmargin=*,label=(\roman*)]

  \item We formulate sensitive microservice placement as a legality-first
  online scheduling problem that distinguishes mandatory policy feasibility
  from temporal operational risk, selection fairness, and commit-time
  resource executability.

  \item We design \cage as an integrated decision pipeline that combines
  deterministic policy gating, calibrated temporal-risk prediction,
  fairness-aware candidate shortlisting, and an empirical Q95 pre-commit
  guardrail. The ordering of these components ensures that prohibited nodes
  never enter performance optimization and that predictive and
  executability checks retain separate roles.

  \item We develop a temporal node-risk model using recent utilization
  levels, trends, ranges, and variability, together with an efficient
  shortlisting mechanism that limits candidate evaluation while controlling
  destination concentration.

  \item We construct a temporally ordered trace-replay evaluation framework
  from the Alibaba Microservices Trace 2021, including strict preprocessing,
  calibration and test separation, natural and service-balanced workloads,
  deterministic policy configurations, sensitivity analysis, ablation
  studies, runtime evaluation, and paired timestamp-cluster statistical
  comparisons.

\end{enumerate}

The remainder of the paper is organized as follows.
Section~\ref{sec:related} reviews microservice scheduling, temporal resource
management, and policy-aware edge orchestration.
Section~\ref{sec:model} introduces the system model and formulates the
scheduling problem.
Section~\ref{sec:method} presents the design and operation of \cage.
Section~\ref{sec:experimental} describes the dataset, replay protocol,
baselines, and experimental settings.
Section~\ref{sec:results} reports the evaluation results.
Section~\ref{sec:discussion} discusses the design implications and practical
trade-offs, and Section~\ref{sec:conclusion} concludes the paper.